\section{Lagrangian 怎样制造全局下界}
\label{sec:09-Convex-Optimization/05-Duality/01}

考虑一个带约束的最小化问题：

\[
\begin{aligned}
\text{minimize}\quad&f_0(x)\\
\text{subject to}\quad&
f_i(x)\le0,\quad i=1,\ldots,m,\\
&h_j(x)=0,\quad j=1,\ldots,p.
\end{aligned}
\]

你还不知道最优解在哪里，却想构造一个对所有可行点都成立的目标下界——不是为了偷懒跳过求解，而是因为任何可靠的下界都能帮你判断一个随手写的候选解离最优还有多远。最理想的情况：不精确求解原问题，也能推出一个有信息量的下限。

约束自身就是原材料。对任意可行点 \(x\)，\(f_i(x)\le0\)，\(h_j(x)=0\)。用非负数 \(\lambda_i\) 乘每个不等式约束，用任意实数 \(\nu_j\) 乘每个等式约束，再加到目标上：

\[
f_0(x)+\sum_{i=1}^m\lambda_i f_i(x)+\sum_{j=1}^p\nu_j h_j(x).
\]

因为 \(\lambda_i\ge0\) 且 \(f_i(x)\le0\)，每个乘积非正，越加越小；等式项乘什么都是零。于是对可行 \(x\)，这个和不大于 \(f_0(x)\)。

由此定义 Lagrangian：

\[
L(x,\lambda,\nu) = f_0(x)
+\sum_{i=1}^m\lambda_i f_i(x)
+\sum_{j=1}^p\nu_j h_j(x),
\qquad \lambda\ge0.
\]

关键不等式：对任何可行 \(x\)，

\[
L(x,\lambda,\nu)\le f_0(x).
\]

乘子 \(\lambda_i\) 可以理解为约束违反量的"罚金单价"——违反得越多，罚金越高；罚金从目标中扣除，而非附加，所以乘子非负才有意义。如果某条约束是松的（\(f_i(x)<0\)），它贡献一个负项，拉低 Lagrangian；如果是紧的（\(f_i(x)=0\)），这一项消失，交叉点处没有惩罚。

\subsection{\texorpdfstring{为什么要对所有 \(x\) 取下确界}{为什么要对所有 x 取下确界}}

你可能想：既然对任意可行 \(x\) 都成立，我就挑一个方便的可行点算个下界。但这样得到的数只保证不大于那一个点的目标值——不能保证它是所有可行点的公共下界。

对所有 \(x\)（不限可行）取下确界：

\[
g(\lambda,\nu) = \inf_x L(x,\lambda,\nu).
\]

现在对于任何可行 \(x\)，

\[
g(\lambda,\nu)\le L(x,\lambda,\nu)\le f_0(x).
\]

下确界不大于 Lagrangian 的每个取值，因此也不大于每个可行目标值，最终不大于原问题最优值 \(p^\star\)——即可行目标集合的下确界：

\[
g(\lambda,\nu)\le p^\star.
\]

每组合法的乘子都免费送你一个 \(p^\star\) 的下界。拿下确界这一步相当于让最不利的 \(x\) 来检验：这组乘子能保证的下界到底有多高。如果 Lagrangian 沿某个方向无界下降，\(g\) 就是 \(-\infty\)——这组乘子没有提供有用信息，但逻辑上它仍然是一个（平凡的）下界。

这里有一个微妙点：下确界是对所有 \(x\) 取的，包括不可行点。在不可行点上，\(\lambda_i f_i(x)\) 可能为正——约束在不可行点上是正的，乘上非负乘子为正，可能让 Lagrangian 远高于原目标。但下确界挑的是最小值，它会自然避开这些"高价"的区域。所以不可行点的存在不影响下界推导的有效性。

\subsection{对偶函数总是凹的}

对偶函数 \(g(\lambda,\nu)\) 是凹函数，即使原始问题非凸。这个性质来自一个极简洁的结构性事实。

对固定 \(x\)，\(L(x,\lambda,\nu)\) 关于 \((\lambda,\nu)\) 是仿射函数——乘子只以线性项出现。任意族仿射函数的逐点下确界是凹函数。验证：对 \(t\in[0,1]\) 与两组乘子 \((\lambda,\nu),(\lambda',\nu')\)，

\[
\begin{aligned}
g\big(t(\lambda,\nu)+(1-t)(\lambda',\nu')\big)
&=\inf_x\Bigl(t\,L(x,\lambda,\nu)+(1-t)\,L(x,\lambda',\nu')\Bigr)\\
&\ge t\,g(\lambda,\nu)+(1-t)\,g(\lambda',\nu'),
\end{aligned}
\]

最后一步因为"和的下确界不小于各自下确界之和"。不等号方向正是凹性所需要的。

既然每组乘子都给下界，自然要寻找最高的一个：

\[
\begin{aligned}
\text{maximize}\quad&g(\lambda,\nu)\\
\text{subject to}\quad&\lambda\ge0.
\end{aligned}
\]

这就是 Lagrange 对偶问题。其最优值记作 \(d^\star\)。注意对偶问题永远是凸的——最大化凹函数，约束 \(\lambda\ge0\) 是半空间。即使原问题非凸，对偶问题也是凸的，这是对偶理论最实用的设计之一。

\subsection{一个一维例子看清全部符号}

\[
\begin{aligned}
\text{minimize}\quad&x^2\\
\text{subject to}\quad&x\ge1.
\end{aligned}
\]

先把约束写成标准形式 \(1-x\le0\)。

Lagrangian：

\[
L(x,\lambda)=x^2+\lambda(1-x),
\qquad\lambda\ge0.
\]

对 \(x\) 取下确界。配方或令导数为零：\(\partial L/\partial x=2x-\lambda=0\) 给出 \(x=\lambda/2\)。代入：

\[
g(\lambda)=\inf_x L(x,\lambda)=\lambda-\frac{\lambda^2}{4}.
\]

开口向下的抛物线。对 \(\lambda\ge0\) 最大化：导数 \(1-\lambda/2=0\) 得 \(\lambda^\star=2\)。所以

\[
d^\star=g(2)=1.
\]

原问题显然在 \(x^\star=1\) 处取得 \(p^\star=1\)。最好的下界恰好碰到原最优值——相差为零。

回头解读 \(\lambda^\star=2\)。在最优解 \(x^\star=1\) 处，约束 \(x\ge1\) 是紧的。如果把约束放松成 \(x\ge0.9\)，最优值会从 \(1\) 降到 \(0.81\)；\(\lambda^\star=2\) 恰好度量了这个边际变化率——约束每放松一单位，最优值改善多少。这是乘子作为"影子价格"的经济直觉，在后续的灵敏度分析中会系统化。

注意乘子必须非负。如果误取 \(\lambda<0\)，对可行点 \(x\ge1\)，\(\lambda(1-x)\ge0\)，于是 \(L(x,\lambda)\ge f_0(x)\)，Lagrangian 不再提供下界。乘子符号不是可选项——它直接保证推导的每一步都有效。

\subsection{弱对偶永远成立}

把所有合法下界取最大，仍不可能超过原最优值：

\[
d^\star\le p^\star.
\]

弱对偶。它不要求原问题凸，不要求最优解存在——只要 Lagrangian 的符号约定正确，结论就无条件成立。

弱对偶立即给出两个实用工具：

\begin{enumerate}
\item
  若找到了原始可行点 \(\hat x\) 和对偶可行乘子 \((\hat\lambda,\hat\nu)\)，则

  \[
  g(\hat\lambda,\hat\nu)
  \le p^\star
  \le f_0(\hat x).
  \]

  上下界之间的间隙 \(f_0(\hat x)-g(\hat\lambda,\hat\nu)\) 告诉你当前的候选解离最优还有多远——在停止优化之前，至少知道不确定性有多大。这个间隙也就是所谓对偶间隙，是终止条件设计的核心依据。

\item
  若某个对偶可行序列的目标趋于 \(+\infty\)，那么原问题不可能有有限可行值——否则弱对偶不等式 \(d^\star\le p^\star\) 会被违反。对偶给出了原问题不可行的证书：你不需要在原空间苦找一个可能不存在的可行点，对偶空间直接宣告了不可行性。
\end{enumerate}

对偶不是从外部强行引入的抽象工具。它从原问题的约束中提取出"这些约束在压低目标"的结构信息，乘子就是量化这个压低效果的参数。下一节要问的是：什么时候这个下界真正 tight——对偶间隙为零，强对偶成立。
